% GENERATED FILE. Edit canonical lessons, then run npm run generate:books.
% canonical-source-hash: fnv1a64:9906ec6d38908d1c
% canonical-lessons: JA-C401-choushoku, JA-C401-yuushoku, JA-C401-chuushoku, JA-C401-kigaeru, JA-C401-motenasu, JA-R401-review-396-401

\chapter{Meals at the Inn}
\label{ch:ja-meals-at-the-inn}
\hlchaptermodality{401}

\begin{chapteropening}
\textbf{By the end of this chapter:} \emph{I can ask about meals at an inn.}

Complete the last lesson of chapter 401: I can ask about meals at an inn.
\end{chapteropening}

\section[\texorpdfstring{\ja{ちょうしょく} (chōshoku)}{ちょうしょく (chōshoku)}]{\ja{ちょうしょく} (chōshoku) — breakfast (a formal word)}

\label{lesson:JA-C401-choushoku}



\begin{quote}
Before the new one: say the Japanese for a toilet, then the Japanese for to shut.
\end{quote}

\subsection*{You'll want to know: \ja{ちょうしょく}}
\textbf{\ja{ちょうしょく}} — \emph{chōshoku} — \textquotedblleft{}breakfast (a formal word)\textquotedblright{}.

\begin{quote}
\textbf{\ja{ちょうしょくは} \ja{なんじからですか。}} — \emph{chōshoku wa nanji kara desu ka} — \textquotedblleft{}What time is breakfast served from?\textquotedblright{}
\end{quote}

\begin{cousinweb}
\textbf{\ja{ちょうしょく}} (\emph{chōshoku}) is Sino-Japanese (kango): two Chinese characters read with their on'yomi, readings borrowed from Chinese. They are \emph{chō}, \textquotedblleft{}morning\textquotedblright{} (Middle Chinese \emph{trjew}) and \emph{shoku}, \textquotedblleft{}eat\textquotedblright{} (\emph{zyik}) (KANJIDIC2; Middle Chinese in Baxter's transcription).
\end{cousinweb}

\begin{grammarlens}[title={What you've built}]
One more word for getting around.
\end{grammarlens}

\subsection*{Your turn}
\begin{itemize}
\raggedright
  \item \emph{Say it:} \emph{chōshoku}
  \item \emph{Say it:} \emph{chōshoku}, once more
  \item \emph{Say it:} \emph{chōshoku wa nanji kara desu ka} — What time is breakfast served from?
  \item \emph{Recall:} say the Japanese for a toilet, then the Japanese for to shut, then say \emph{chōshoku} again
\end{itemize}

\begin{tcolorbox}[breakable,colback=teal!4,colframe=teal!35!black,title={Before you move on}]
What does \ja{ちょうしょく} mean? (\textquotedblleft{}Breakfast (a formal word)\textquotedblright{}.) Say it once more.
\end{tcolorbox}
\section[\texorpdfstring{\ja{ゆうしょく} (yūshoku)}{ゆうしょく (yūshoku)}]{\ja{ゆうしょく} (yūshoku) — dinner (a formal word)}

\label{lesson:JA-C401-yuushoku}



\begin{quote}
Before the new one: say the Japanese for to shut, then the Japanese for breakfast (a formal word).
\end{quote}

\subsection*{You'll want to know: \ja{ゆうしょく}}
\textbf{\ja{ゆうしょく}} — \emph{yūshoku} — \textquotedblleft{}dinner (a formal word)\textquotedblright{}.

\begin{quote}
\textbf{\ja{りょかんの} \ja{ゆうしょくは} \ja{さかなでした。}} — \emph{ryokan no yūshoku wa sakana deshita} — \textquotedblleft{}Dinner at the inn was fish.\textquotedblright{}
\end{quote}

\begin{cousinweb}
\textbf{\ja{ゆうしょく}} (\emph{yūshoku}) is built from both layers of the vocabulary: the native Japanese (wago) \emph{yū}, \textquotedblleft{}evening\textquotedblright{}, inherited rather than borrowed, and the Sino-Japanese \emph{shoku}, \textquotedblleft{}eat\textquotedblright{} (Middle Chinese \emph{zyik}), an on'yomi borrowed from Chinese (JMdict; KANJIDIC2; Middle Chinese in Baxter's transcription).
\end{cousinweb}

\begin{grammarlens}[title={What you've built}]
One more word for getting around.
\end{grammarlens}

\subsection*{Your turn}
\begin{itemize}
\raggedright
  \item \emph{Say it:} \emph{yūshoku}
  \item \emph{Say it:} \emph{yūshoku}, once more
  \item \emph{Say it:} \emph{ryokan no yūshoku wa sakana deshita} — Dinner at the inn was fish.
  \item \emph{Recall:} say the Japanese for to shut, then the Japanese for breakfast (a formal word), then say \emph{yūshoku} again
\end{itemize}

\begin{tcolorbox}[breakable,colback=teal!4,colframe=teal!35!black,title={Before you move on}]
What does \ja{ゆうしょく} mean? (\textquotedblleft{}Dinner (a formal word)\textquotedblright{}.) Say it once more.
\end{tcolorbox}
\section[\texorpdfstring{\ja{ちゅうしょく} (chūshoku)}{ちゅうしょく (chūshoku)}]{\ja{ちゅうしょく} (chūshoku) — lunch (a formal word)}

\label{lesson:JA-C401-chuushoku}



\begin{quote}
Before the new one: say the Japanese for breakfast (a formal word), then the Japanese for dinner (a formal word).
\end{quote}

\subsection*{You'll want to know: \ja{ちゅうしょく}}
\textbf{\ja{ちゅうしょく}} — \emph{chūshoku} — \textquotedblleft{}lunch (a formal word)\textquotedblright{}.

\begin{quote}
\textbf{\ja{ちゅうしょくは} \ja{えきべんでした。}} — \emph{chūshoku wa ekiben deshita} — \textquotedblleft{}Lunch was a station lunch box.\textquotedblright{}
\end{quote}

\begin{cousinweb}
\textbf{\ja{ちゅうしょく}} (\emph{chūshoku}) is Sino-Japanese (kango): two Chinese characters read with their on'yomi, readings borrowed from Chinese. They are \emph{chū}, \textquotedblleft{}daytime\textquotedblright{} (Middle Chinese \emph{trjuwH}) and \emph{shoku}, \textquotedblleft{}eat\textquotedblright{} (\emph{zyik}) (KANJIDIC2; Middle Chinese in Baxter's transcription).
\end{cousinweb}

\begin{grammarlens}[title={What you've built}]
One more word for getting around.
\end{grammarlens}

\subsection*{Your turn}
\begin{itemize}
\raggedright
  \item \emph{Say it:} \emph{chūshoku}
  \item \emph{Say it:} \emph{chūshoku}, once more
  \item \emph{Say it:} \emph{chūshoku wa ekiben deshita} — Lunch was a station lunch box.
  \item \emph{Recall:} say the Japanese for breakfast (a formal word), then the Japanese for dinner (a formal word), then say \emph{chūshoku} again
\end{itemize}

\begin{tcolorbox}[breakable,colback=teal!4,colframe=teal!35!black,title={Before you move on}]
What does \ja{ちゅうしょく} mean? (\textquotedblleft{}Lunch (a formal word)\textquotedblright{}.) Say it once more.
\end{tcolorbox}
\section[\texorpdfstring{\ja{きがえる} (kigaeru)}{きがえる (kigaeru)}]{\ja{きがえる} (kigaeru) — to change one's clothes}

\label{lesson:JA-C401-kigaeru}



\begin{quote}
Before the new one: say the Japanese for dinner (a formal word), then the Japanese for lunch (a formal word).
\end{quote}

\subsection*{You'll want to know: \ja{きがえる}}
\textbf{\ja{きがえる}} — \emph{kigaeru} — \textquotedblleft{}to change one's clothes\textquotedblright{}.

\begin{quote}
\textbf{\ja{おふろの} \ja{あとで} \ja{ゆかたに} \ja{きがえました。}} — \emph{ofuro no ato de yukata ni kigaemashita} — \textquotedblleft{}After the bath I changed into a cotton kimono.\textquotedblright{}
\end{quote}

Polite \textbf{\ja{きがえます}}, \emph{kigaemasu}; an ichidan verb. Past (polite): \textbf{\ja{あにが} \ja{ゆかたに} \ja{きがえました。}}, \emph{ani ga yukata ni kigaemashita}.

\begin{cousinweb}
\textbf{\ja{きがえる}} (\emph{kigaeru}) is a native Japanese word (wago), inherited rather than borrowed from Chinese, built from \emph{ki}, \textquotedblleft{}wearing\textquotedblright{} (from \emph{kiru}) and \emph{gaeru}, \textquotedblleft{}changing\textquotedblright{} (from \emph{kaeru}); its kanji were attached later as kun'yomi (JMdict).
\end{cousinweb}

\begin{grammarlens}[title={What you've built}]
One more word for getting around.
\end{grammarlens}

\subsection*{Your turn}
\begin{itemize}
\raggedright
  \item \emph{Say it:} \emph{kigaeru}
  \item \emph{Say it:} \emph{kigaeru}, once more
  \item \emph{Say it:} \emph{ofuro no ato de yukata ni kigaemashita} — After the bath I changed into a cotton kimono.
  \item \emph{Recall:} say the Japanese for dinner (a formal word), then the Japanese for lunch (a formal word), then say \emph{kigaeru} again
\end{itemize}

\begin{tcolorbox}[breakable,colback=teal!4,colframe=teal!35!black,title={Before you move on}]
What does \ja{きがえる} mean? (\textquotedblleft{}To change one's clothes\textquotedblright{}.) Say it once more.
\end{tcolorbox}
\section[\texorpdfstring{\ja{もてなす} (motenasu)}{もてなす (motenasu)}]{\ja{もてなす} (motenasu) — to entertain guests; to look after visitors}

\label{lesson:JA-C401-motenasu}



\begin{quote}
Before the new one: say the Japanese for lunch (a formal word), then the Japanese for to change one's clothes.
\end{quote}

\subsection*{You'll want to know: \ja{もてなす}}
\textbf{\ja{もてなす}} — \emph{motenasu} — \textquotedblleft{}to entertain guests; to look after visitors\textquotedblright{}.

\begin{quote}
\textbf{\ja{りょかんは} \ja{きゃくを} \ja{ていねいに} \ja{もてなしました。}} — \emph{ryokan wa kyaku o teinei ni motenashimashita} — \textquotedblleft{}The inn looked after its guests with great care.\textquotedblright{}
\end{quote}

Polite \textbf{\ja{もてなします}}, \emph{motenashimasu}; a godan verb. Past (polite): \textbf{\ja{そぼが} \ja{きゃくを} \ja{もてなしました。}}, \emph{sobo ga kyaku o motenashimashita}.

\begin{cousinweb}
\textbf{\ja{もてなす}} (\emph{motenasu}) is a native Japanese word (wago), inherited rather than borrowed from Chinese (JMdict).
\end{cousinweb}

\begin{grammarlens}[title={What you've built}]
One more word for getting around.
\end{grammarlens}

\subsection*{Your turn}
\begin{itemize}
\raggedright
  \item \emph{Say it:} \emph{motenasu}
  \item \emph{Say it:} \emph{motenasu}, once more
  \item \emph{Say it:} \emph{ryokan wa kyaku o teinei ni motenashimashita} — The inn looked after its guests with great care.
  \item \emph{Recall:} say the Japanese for lunch (a formal word), then the Japanese for to change one's clothes, then say \emph{motenasu} again
\end{itemize}

\begin{tcolorbox}[breakable,colback=teal!4,colframe=teal!35!black,title={Before you move on}]
What does \ja{もてなす} mean? (\textquotedblleft{}To entertain guests; to look after visitors\textquotedblright{}.) Say it once more.
\end{tcolorbox}
\section[(dialogue)]{Looking back: from Where Are You Headed? to Meals at the Inn}

\label{lesson:JA-R401-review-396-401}



\begin{quote}
Say the Japanese for a destination and to wander about.
\end{quote}

\subsection*{Your turn}
Say these aloud:

\begin{itemize}
\raggedright
  \item a destination, to wander about, to finally reach, strolling, wandering restlessly
  \item a family-run guesthouse, accommodation, a one-night stay, a day trip, to look after
  \item a hot spring, rising steam, to feel dizzy, to soak, to relax
  \item a pillow, a blanket, a floor cushion, to wake from sleep, to lie down
  \item an alcove, a tatami room, a veranda, a toilet, to shut
  \item breakfast (a formal word), dinner (a formal word), lunch (a formal word), to change one's clothes, to entertain guests
\end{itemize}

\begin{tcolorbox}[breakable,colback=teal!4,colframe=teal!35!black,title={Before you move on}]
A destination? (\textbf{\ja{ゆきさき}}, \emph{yukisaki}.) To wander about? (\textbf{\ja{さまよう}}, \emph{samayou}.) To finally reach? (\textbf{\ja{たどりつく}}, \emph{tadoritsuku}.) Strolling? (\textbf{\ja{ぶらぶら}}, \emph{burabura}.) Wandering restlessly? (\textbf{\ja{うろうろ}}, \emph{urōro}.) A family-run guesthouse? (\textbf{\ja{みんしゅく}}, \emph{minshuku}.) Accommodation? (\textbf{\ja{しゅくはく}}, \emph{shukuhaku}.) A one-night stay? (\textbf{\ja{いっぱく}}, \emph{ippaku}.) A day trip? (\textbf{\ja{ひがえり}}, \emph{higaeri}.) To look after? (\textbf{\ja{あずかる}}, \emph{azukaru}.) A hot spring? (\textbf{\ja{おんせん}}, \emph{onsen}.) Rising steam? (\textbf{\ja{ゆげ}}, \emph{yuge}.) To feel dizzy? (\textbf{\ja{のぼせる}}, \emph{noboseru}.) To soak? (\textbf{\ja{つかる}}, \emph{tsukaru}.) To relax? (\textbf{\ja{くつろぐ}}, \emph{kutsurogu}.) A pillow? (\textbf{\ja{まくら}}, \emph{makura}.) A blanket? (\textbf{\ja{もうふ}}, \emph{mōfu}.) A floor cushion? (\textbf{\ja{ざぶとん}}, \emph{zabuton}.) To wake from sleep? (\textbf{\ja{めざめる}}, \emph{mezameru}.) To lie down? (\textbf{\ja{ねころぶ}}, \emph{nekorobu}.) An alcove? (\textbf{\ja{とこのま}}, \emph{tokonoma}.) A tatami room? (\textbf{\ja{ざしき}}, \emph{zashiki}.) A veranda? (\textbf{\ja{えんがわ}}, \emph{engawa}.) A toilet? (\textbf{\ja{おてあらい}}, \emph{otearai}.) To shut? (\textbf{\ja{しめる}}, \emph{shimeru}.) Breakfast (a formal word)? (\textbf{\ja{ちょうしょく}}, \emph{chōshoku}.) Dinner (a formal word)? (\textbf{\ja{ゆうしょく}}, \emph{yūshoku}.) Lunch (a formal word)? (\textbf{\ja{ちゅうしょく}}, \emph{chūshoku}.) To change one's clothes? (\textbf{\ja{きがえる}}, \emph{kigaeru}.) To entertain guests? (\textbf{\ja{もてなす}}, \emph{motenasu}.)
\end{tcolorbox}
